\section{C. 정렬의 사회성}

\begin{frame} % No title at first slide
    \sectiontitle{C}{정렬의 사회성}
    \sectionmeta{
        출제자 -- 장영후({\href{https://solved.ac/profile/yhjang04}{@yhjang04}})\\
        태그: \consolas {\#string, \#sorting, \#implementation}\\
        예상 난이도 -- \textbf{\color{acS4}Silver4}
    }
    \begin{itemize}
        \item 제출 49번, 정답 10명 (정답률 20.4\%)
        \item 처음 푼 사람: 김0건 19분
    \end{itemize}
\end{frame}
%----------------------------------------------------------------------------
\begin{frame}{\textbf{C}. 정렬의 사회성}
    \begin{itemize}
        \item 가장 중요한 점은 두 단어의 대소비교 함수를 구현하는 것입니다.\\
        문자별 인덱스를 기록하여, 이 인덱스로 비교하는 것이 정해입니다.

        \item 다만 문자열의 길이 비교까지도 통상적인 정렬 방식과 다름을 유의하여야 합니다.

        \item 직접 정렬을 구현해서 사용하거나, 혹은 각 언어별로 지원하는 정렬 함수를 사용하여도 됩니다.\\
        단어의 개수가 1,000개 이하이므로 대부분의 경우 $O(n^2)$으로 정렬을 구현하더라도 통과할 수 있습니다.
    \end{itemize}
\end{frame}